\title{Frozen Basis Time Series Forecasting}

\author{%
  Daniel Byrne\\
  \texttt{realdanielbyrne@icloud.com}
}

\maketitle

\begin{abstract}
N-BEATS reaches its M4 accuracy with twenty to forty five million parameters, a capacity never
questioned. Which parts of a basis expansion block carry the accuracy, and which are inert? I
factor each block's unconstrained synthesis head through a frozen orthonormal basis, leaving two
integers, the backcast and forecast ranks, plus one trunk choice, varying each at matched
budgets.

The parameters are in the trunk and are not needed: narrowing it from 2048 units to 128 cuts the
count 48.8 fold, 64,645,150 to 1,323,550, for 0.098 sMAPE at p = 0.78. The compression credited
to the wavelet basis is the trunk, one integer in the original architecture. Basis identity
cannot matter at the rank every reported configuration uses, where the basis is square and
orthogonal and the projection feeding it absorbs any rotation exactly; thirty six paired
comparisons over wavelet, cosine and random bases are null, minimum p = 0.084. Below full rank
it matters greatly, and captured energy, the share of target energy the retained rows span,
reproduces the trained ordering without inversion on two datasets, rank correlation -1.0 on
Traffic-96. Trainability is a separate axis: effective rank tracks convergence at rho = +0.949,
bases concentrating their energy in one direction fail on every seed, and rectifying the output
or pairing a trend branch restores full convergence without adding expressive power. Instability
is negative forecasts on positive series, which sMAPE scores at its 200 ceiling regardless of
magnitude; deleting those points recovers the converged score to within 0.11, and rectification
cuts collapses from 107 of 23,786 runs to 0 of 136. Out of sample on Tourism the compact
configuration matches or beats both paper baselines at 14 to 15 times fewer parameters, and
never collapses in 240 runs.

The gain is position on the accuracy against size frontier, not accuracy: two significant losses
on M4 and an outright loss on Traffic-96, 29.37 sMAPE at 1.81M against 26.24 at 52.3M.
\end{abstract}

\section{Introduction}\label{sec:intro}

The submission this rewrite replaces claimed that wavelets and autoencoders deliver accuracy and
compression on top of the N-BEATS architecture \citep{oreshkin2020nbeats}; neither half survived
the controls run after submission. The corrected thesis is narrower and, we argue, more useful.
N-BEATS carries roughly fifty times more trunk capacity than its single-model accuracy requires,
where the \emph{trunk} is the stack of fully connected layers inside each block that maps the
input window to a hidden vector. Replacing each block's unconstrained linear synthesis head with
a frozen orthonormal \emph{basis} --- a fixed set of orthonormal vectors into which the head's
output is expanded --- makes the rank an explicit knob and removes the head's ability to emit
values outside the target's range. Which basis fills that slot is provably irrelevant once it is
complete. Two quantities computable before training, captured energy and effective rank, say how
accurate a truncated basis will be and whether it will train at all.

Four results carry the paper. First, the trunk result (Section~\ref{sec:trunk}): trunk width is
inert across a 48.8-fold parameter range --- 160 runs, rung means spanning 0.098~\smape{}
against within-rung seed standard deviations of 0.15--0.25, Kruskal--Wallis $p = 0.776$ and
$0.773$ --- so the compression the submission credited to the basis was one integer in a
constructor. Second, the completeness result (Section~\ref{sec:completeness}): at complete
forecast rank with linear synthesis, basis identity is provably inert --- the learned projection
absorbs any rotation exactly, at residual $5.6 \times 10^{-15}$, and none of 36 paired
comparisons is significant, minimum $p = 0.084$ --- subject to the four scope conditions stated
in Section~\ref{sec:method}. Third, the truncated-rank result (Section~\ref{sec:truncation}):
captured energy, the share of target energy the retained basis rows span, predicts the trained
ordering exactly (Traffic-96, five bases, Spearman $\rho = -1.0$, exact $p = 0.017$; run-level
$\rho = -0.98$, $p = 9.0 \times 10^{-18}$). Fourth, the trainability result
(Section~\ref{sec:trainability}): at complete rank the DCT and Karhunen--Lo\`eve bases fail on
every seed under a plain linear head, with effective rank $1.05$ and $1.54$ against $17.9$ for
bases that never fail, and a paired trend branch repairs both to 10/10 converged seeds
($p = 0.00018$).

Three things this paper does not claim. It does not claim that wavelets beat Fourier: at matched
width that comparison is null. It does not claim wavelet-family effects across datasets: that
claim is retracted, Fisher combined $p = 0.125$. It does not claim anything about Transformers:
the submission's Transformer section is withdrawn (Section~\ref{sec:scoped}). What remains is a
measurement about N-BEATS, together with the design recipe that follows from it.

Section~\ref{sec:method} states the method in two integers and one trunk choice;
Sections~\ref{sec:trunk}--\ref{sec:trainability} develop the four load-bearing results;
Section~\ref{sec:tourism} validates the compact configuration out of sample;
Section~\ref{sec:scoped} collects the scoped, transferred, and withdrawn claims;
Section~\ref{sec:related_recipe_limits} positions the work against prior art, gives the
practitioner's recipe, and states the limitations. Appendices carry the protocols, the full
per-period tables, the configuration lists, the out-of-distribution detail, and the reproduction
material; all configurations, per-run CSVs, and registered prediction files ship with the paper.
